\documentclass[a4paper]{article}
\usepackage[pdftex]{graphicx}
\usepackage[utf8]{inputenc}
\usepackage{enumerate}
\usepackage{siunitx}
\sisetup{locale=DE}
\usepackage{href-ul}
\hypersetup{
	colorlinks=true,
	linkcolor=blue,
	urlcolor=blue}
\usepackage{icomma}
\usepackage{amssymb}
\usepackage{geometry}
\geometry{a4paper, top=15mm, left=15mm, right=15mm, bottom=15mm,
	headsep=10mm, footskip=12mm}
\hyphenation{weight-force}


\begin{document}
	%Title
	{\bf School assignment from physics }\\
	\begin{enumerate}[1.]
		
		%Task 1
		{\bf \item Task}
		
		\begin{enumerate}[a)]
			\item Explain the creation of sliding friction using a sketch.
			\item Name two measures how you can reduce sliding friction.
			\item Explain the difference between one-sided and two-sided levers and give three examples of each type.
			\item What is meant by torque?
			\item What is the law of conservation of energy?
		\end{enumerate}
		
		%Exercise 2
		{\bf \item In order to determine the dependence of the sliding friction force on the contact force, the following series of measurements was recorded:}
		
		\begin{tabular}{ccccccc}
			$F_N$ in N & 0 & 0.5 & 1.5 & 2.0 & 2.5 & 3.0 \\
			$F_{RG}$ in N & 0 & 0.16 & 0.47 & 0.62 & 0.78 & 0.91
		\end{tabular}
		
		\begin{enumerate}[a)]
			\item Display the series of measurements graphically.
			\item Formulate the legal connection that results.
			\item Calculate the coefficient of sliding friction $\mu$ from the diagram.
		\end{enumerate}
		
		
		%Task 5
		{\bf \item Three equally sized cuboids are to be pulled across a wooden table.} Think about the possible arrangement options and calculate the pulling force to be applied for the best option.\\
		
		\renewcommand{\arraystretch}{1.5}
		\begin{tabular}{ccccc}
			\framebox[2.2 cm]{Wood} & $F_G=\SI{5.0}{\newton}$ & $\mu_{Wood-Wood}$ &=& 0.30 \\
			\framebox[2.2 cm]{Iron} & $F_G=\SI{15}{\newton}$ & $\mu_{Iron-Wood}$ &=& 0.40 \\
			\framebox[2.2 cm]{plastic} & $F_G=\SI{3.0}{\newton}$ & $\mu_{plastic-wood}$ &=& 0.35
		\end{tabular}
		
		%Task 6
		{\bf \item To what height was a pallet with beverage crates with the weight of $F_G=\SI{1500}{\newton}$
			lifted when a work of $W=\SI{12}{\kilo\joule}$ was done on it?}
		
		%Task 7
		{\bf \item A cabinet is moved in $\SI{10}{\second}$ by $\SI{2.0}{\meter}$ ($\mu = 0.20 $).} For this is one to generate power of $P=\SI{20}{\watt}$. What is the weight of the cabinet?
		
	\end{enumerate}
	\fbox{
		\begin{minipage}{0.5\textwidth}
			For the solution, please \href{https://www.okuyakl.de/phys/p8reibL122/le122.pdf}{click here} or scan the QR code.\\
			Additional worksheets are available at
			
			\href{http://www.okuyakl.com}{www.okuyakl.com}
		\end{minipage}
		\hfill
		\begin{minipage}{0.4\textwidth}
			\includegraphics[width=1.5 cm]{../../viecher/zwe03}
			\includegraphics[width=3 cm]{qre122}
			\includegraphics[width=2 cm]{../../viecher/afanticon1}
			
	\end{minipage}}
\end{document}